% !TEX root = ../DoAn.tex
\documentclass[../DoAn.tex]{subfiles}
\begin{document}

\begin{center}
    \Large{\textbf{TÓM TẮT NỘI DUNG ĐỒ ÁN}}\\
\end{center}
\vspace{0.5cm}

Hệ thống hỏi--đáp dựa trên truy hồi (Retrieval-Augmented Generation, RAG) cho tài liệu kỹ thuật tiếng Việt đối mặt với hai thách thức: chất lượng truy hồi trên văn bản chuyên ngành nhiều thuật ngữ và bảng biểu, cùng chi phí và độ trễ của tầng xếp hạng lại. Các giải pháp hiện tại thường dùng cross-encoder lớn như BGE-reranker-v2-m3 (568M tham số) để xếp hạng lại và một mô hình ngôn ngữ lớn qua API để định tuyến câu hỏi --- cho độ chính xác cao nhưng nặng, chậm trên CPU và phát sinh chi phí, trong khi embedding nền tảng chưa được tối ưu cho miền dữ liệu. Đồ án tối ưu toàn bộ pipeline thành các thành phần nhẹ, chạy được trên CPU với chi phí thấp mà vẫn giữ chất lượng gần với mô hình lớn.

Embedding \plaincode{Vietnamese_Embedding_v2} được tinh chỉnh bằng hàm mất mát MultipleNegativesRankingLoss trên các cặp câu hỏi--positive của miền, kết hợp Matryoshka Representation Learning để một mô hình duy nhất cho biểu diễn dùng được ở ba số chiều 256/512/1024; cấu hình 512 chiều được chọn để triển khai, giảm một nửa dung lượng vector index. Reranker được chưng cất tri thức từ teacher BGE-reranker-v2-m3 sang student mmarco-mMiniLMv2 117M và một bản cắt tỉa từ vựng 35.6M bằng hàm mất mát ADR-MSE dựa trên thứ hạng, rồi lượng tử hoá INT8 để chạy trên CPU. Bộ định tuyến ý định dùng SetFit trên nền Sentence-BERT tiếng Việt, thay thế hoàn toàn lời gọi API.

Kết quả (trung bình trên ba hạt giống ngẫu nhiên), embedding tinh chỉnh cải thiện Acc@1 từ 65.19\% lên $73.21 \pm 0.60$\% trên tập in-domain và đạt MRR@10 $= 0.7982 \pm 0.0017$ ở 512 chiều trên tập tài liệu chưa từng xuất hiện khi huấn luyện. Reranker 35.6M đạt MRR@10 $= 0.8844$, vượt ViRanker 568M (0.8593) dù nhỏ hơn khoảng 16 lần; student 117M đạt $0.8862$, bằng đúng teacher 568M. Sau lượng tử hoá INT8, reranker chỉ còn 34.8 MB và nhanh hơn teacher khoảng 10 lần trên CPU. Đánh giá đầu cuối trên 390 câu trắc nghiệm cho thấy pipeline tối ưu giữ nguyên độ chính xác của hệ thống ban đầu (97.4\%) trong khi phần truy hồi--xếp hạng nhanh hơn khoảng 20 lần, chứng minh tính khả thi của một hệ thống RAG tiếng Việt hiệu quả với chi phí thấp.

\begin{flushright}
\begin{tabular}{c}
Sinh viên thực hiện\\[0.2cm]
\textit{Huy}\\
Trần Quang Huy
\end{tabular}
\end{flushright}

\end{document}